\chapter{Base-Direction Information Geometry}
\label{ch:pullback-geometry}

The fiber geometry of the main manuscript lives on one associated fiber: the Fisher metric
\(g^F_x\) and the Amari--Chentsov tensor \(T^F_x\) of the regular channel models
(\Cref{hyp:pb-regular-models}), their descent to the associated bundles
(\Cref{prop:pb-statistical-tensor-descent}), and the vertical Fisher defect of a Markov coarse
morphism with its cocycle (\Cref{thm:pb-pullback-fisher-defect,thm:pb-fisher-defect-cocycle}). Over
\(\{c_*\}\) the base tangent space is zero, so the covariant first jet \(D^\omega s\) of a section,
its pullback tensors \(h_s^\omega\) and \(c_s^\omega\), the horizontal defect \(A_\Psi\) of a
bundle morphism, and the isotropy criterion below all vanish identically. This chapter asks what the
fiber geometry induces on the base once contexts vary and connections are chosen. The pullbacks are
gauge invariant but connection relative, with exact change formulas
(\Cref{thm:pb-pullback-gauge-invariance,prop:pb-pullback-connection-change}). Across a bundle
morphism over a base map, fine and coarse pullbacks differ by an exact signed identity with two
retained anomaly tensors (\Cref{thm:pb-signed-base-comparison}); the horizontal defect of an induced
morphism is the fundamental field of a scale-connection defect form and vanishes exactly under an
isotropy criterion (\Cref{thm:pb-isotropy-criterion}); and the base Fisher defect composes by a
cocycle up to an explicit residual (\Cref{thm:pb-base-defect-cocycle}), whose relaxed pointwise
defect problems have complete finite-dimensional parameterizations
(\Cref{thm:pb-defect-contractions,thm:pb-defect-partial-isometries,thm:pb-defect-signed-gram}). These
results are used by the section-space histories of \Cref{ch:section-space-histories}, the covers of
\Cref{ch:neighborhood-parents}, and the perceived-geometry beta of \Cref{ch:renormalization}. At
\(c_*\) the counterpart of a scale-connection defect is the defect of the dressed interaction links
under agent coarse-graining, treated in \Cref{ch:network-state} of the main manuscript.

No connection-independent tensor on the base is obtained from the bundle and the statistical models
alone, no canonical connection is derived from the change formulas or from passive covariance, and
the pullback tensors are not claimed to be nondegenerate. A scalar gauged sigma energy would further
require a base cometric, a base density, channel weights, boundary conditions, and a decision whether
the connection is fixed or dynamical, none of which is selected here. The classifications of
\Cref{sec:pb-intrinsic-defect-classification} assert no global solution with prescribed boundary data
across changing ranks, no infinite-dimensional closed-range result, and no compatible connection for
every scale map. \status{NOT-CLAIMED}

\section{Covariant first jets and pullback tensors}
\label{sec:pb-covariant-jets}

Fix one channel \(x\) and write \(E=\mathcal E_x\), \(\mathcal B=\mathcal B_x\),
\(\varpi=\varpi_x\), \(g^F=g^F_x\), \(T^F=T^F_x\), and \(\omega=\omega_x\), with the fiber tensors of
\cref{eq:pb-fiber-fisher,eq:pb-fiber-amari} descended as in \Cref{prop:pb-statistical-tensor-descent}.
The connection induces a horizontal subbundle \(H^\omega E\subset TE\) and a vertical projection
\(\operatorname{ver}^\omega:TE\to VE\). For \(e\in E\) write
\(\mathsf J_e:\operatorname{Ad}(\mathscr P_G)_{\varpi(e)}\to V_eE\) for the fundamental vertical map
of the represented action, and \(\mathsf f_\xi\) for the fundamental vertical field of
\(\xi\in\mathfrak g\) in a local frame.

\definitionheading{Covariant first jets and informational pullbacks}{def:pb-informational-pullbacks}
For a smooth section \(s:\mathcal U\to E\) over an open set \(\mathcal U\subseteq\mathcal C\), the
covariant first jet is
\begin{equation}
D^\omega s
:=\operatorname{ver}^{\omega}\circ Ts
:T\mathcal U\longrightarrow s^*VE,
\label{eq:pb-covariant-first-jet}
\end{equation}
and the connection-relative Fisher and Amari--Chentsov pullbacks of \(s\) are
\begin{align}
h_s^\omega(X,Y)
&=g^F_{s(c)}(D^\omega s(X),D^\omega s(Y)),
\label{eq:pb-fisher-pullback}\\
c_s^\omega(X,Y,Z)
&=T^F_{s(c)}(D^\omega s(X),D^\omega s(Y),D^\omega s(Z)),
\label{eq:pb-amari-pullback}
\end{align}
for \(X,Y,Z\in T_c\mathcal C\). Thus \(h_s^\omega\in\Gamma(\operatorname{Sym}^2T^*\mathcal U)\) and
\(c_s^\omega\in\Gamma(\operatorname{Sym}^3T^*\mathcal U)\); the first is called a semimetric until
its nondegeneracy is proved.

For nonlinear statistical fibers, \(D^\omega s\) is a vertical-valued one-form along \(s\), not a
linear covariant derivative on a vector space of sections. In the Gaussian realization of
Supplement G the fiber metric is the moment-chart Fisher metric of \Cref{prop:ig-fisher-moment-chart},
and \Cref{def:pb-informational-pullbacks} pulls it back to a base tensor that depends on the section
and the connection; by \Cref{thm:pb-pullback-rank-quotient} its rank is the rank of the first jet.
The Gaussian component formulas do not identify the base with the sample or parameter spaces.

\theoremheading{Passive gauge covariance of the covariant pullbacks}{thm:pb-pullback-gauge-invariance}
Under a passive local change of frame \(u'=u\cdot g^{-1}\), with the connection and the section
coordinates transformed by their associated laws, \(h_s^\omega\) and \(c_s^\omega\) are unchanged as
tensors on \(\mathcal U\). They are not invariant under active gauge transformations at a fixed
connection: on the trivial principal \(\R\)-bundle over \(\R\) with the normal location fiber
\(\{\mathcal N(\mu,1)\}\), the zero connection, and the zero section, \(h=0\), while the bundle
automorphism \(F(u(x)g)=u(x)(x+g)\) carries that section to \(s'(x)=\mathcal N(x,1)\) and \(h\) to
\(dx^2\). \status{ESTABLISHED}

\noindent\emph{Proof.}
In a local frame the vertical jet transforms by the tangent representation,
\begin{equation}
D^{\omega'}s'(X)
=T\widehat\rho_x(g)D^\omega s(X),
\label{eq:pb-covariant-jet-gauge-law}
\end{equation}
because coordinates change by \(\beta'=\widehat\rho_x(g)\beta\) under the quotient convention
\eqref{eq:geo-quotient-convention}, and the connection changes by \eqref{eq:geo-local-connection-b}
or \eqref{eq:geo-local-connection-m}. By \Cref{prop:pb-statistical-tensor-descent} the tangent
representation is an isometry of \(g^F\) and preserves \(T^F\), so the frame change cancels in every
argument of \eqref{eq:pb-fisher-pullback} and \eqref{eq:pb-amari-pullback}. For the active example,
with zero connection the vertical jet of a section of means \(\mu(x)\) is \(d\mu\), and the location
Fisher metric is \(1\). \hfill$\square$

Passive covariance says that two coordinate descriptions of one pair \((\omega,s)\) give the same
tensor. It does not say that two connections give the same pullback, which is the subject of the
next section.

\section{Exact dependence on the chosen connection}
\label{sec:pb-connection-dependence}

\definitionheading{Connection difference and its vertical image}{def:pb-connection-difference}
For a second principal connection \(\omega'\) on \(\mathscr P_G\), the difference \(\omega'-\omega\)
is horizontal and \(\operatorname{Ad}\)-equivariant; its descended form is
\begin{equation}
\mathfrak d\in\Omega^1(\mathcal C,\operatorname{Ad}(\mathscr P_G)).
\label{eq:pb-connection-difference}
\end{equation}
Along a section \(s\), its vertical image is
\begin{equation}
\mathsf J^s_{\mathfrak d}(X)=\mathsf J_{s(c)}(\mathfrak d_c(X)).
\label{eq:pb-connection-difference-vertical}
\end{equation}
The sign convention is that in a local trivialization the horizontal lift of \(X\) at \((c,\beta)\) is
\((X,-\mathsf f_{A(X)}\beta)\), compatibly with \eqref{eq:geo-quotient-convention}. For a law fiber that
is not a vector space the linear expression \(\widehat\rho_{x*}(A)\) is undefined; it is available only
in a fiberwise-linear realization, where it agrees with \(\mathsf f_{A(X)}\).

\propositionheading{Connection-change formulas}{prop:pb-pullback-connection-change}
For \(\omega'=\omega+\mathfrak d\), the vertical jets and Fisher pullbacks obey
\begin{align}
D^{\omega'}s
&=D^\omega s+\mathsf J^s_{\mathfrak d},
\label{eq:pb-jet-connection-change}\\
h_s^{\omega'}(X,Y)
&=h_s^\omega(X,Y)
+g^F(\mathsf J^s_{\mathfrak d}X,D^\omega sY)
+g^F(D^\omega sX,\mathsf J^s_{\mathfrak d}Y)
+g^F(\mathsf J^s_{\mathfrak d}X,\mathsf J^s_{\mathfrak d}Y),
\label{eq:pb-fisher-connection-change}
\end{align}
and the Amari--Chentsov pullback satisfies the polarized identity
\begin{equation}
c_s^{\omega'}(X,Y,Z)
=T^F\bigl((D^\omega s+\mathsf J^s_{\mathfrak d})X,
(D^\omega s+\mathsf J^s_{\mathfrak d})Y,
(D^\omega s+\mathsf J^s_{\mathfrak d})Z\bigr),
\label{eq:pb-amari-connection-change}
\end{equation}
which contains the original term, three terms with one occurrence of \(\mathsf J^s_{\mathfrak d}\),
three with two, and one with three. The dependence occurs already for the trivial translation bundle
over \(\R\) with fiber \(\{\mathcal N(\mu,1)\}\) and the constant section \(\mu=0\): the zero
connection gives \(h=0\), and the local connection \(A'=\mathrm c_1dx\) gives
\begin{equation}
h^{A'}=\mathrm c_1^2dx^2.
\label{eq:pb-connection-dependence-example}
\end{equation}
\status{ESTABLISHED}

\noindent\emph{Proof.}
In one local frame, replacing the horizontal lift \((X,-\mathsf f_{A(X)}\beta)\) by
\((X,-\mathsf f_{A(X)+\mathfrak d(X)}\beta)\) shifts the vertical part of \(Ts\) by
\(\mathsf J^s_{\mathfrak d}(X)\), which proves \eqref{eq:pb-jet-connection-change}. Bilinearity of
\(g^F\) and trilinearity of \(T^F\) give the other two formulas, which are frame independent because
\(\mathfrak d\) is \(\operatorname{Ad}(\mathscr P_G)\)-valued and every contraction is gauge
invariant. In the example the vertical jet of the constant section is \(A'(X)\partial_\mu\) and the
Fisher metric is \(1\). \hfill$\square$

The example refutes a connection-independent reading of \(h_s^\omega\); it does not refute the gauge
invariance of a fixed connection and section.

\section{Two channel geometries}
\label{sec:pb-two-channels}

Restoring the channel labels gives the ordered pair
\begin{equation}
(h_{i,b}^{\omega_b},h_{i,m}^{\omega_m})
\in\Gamma(\operatorname{Sym}^2T^*\mathcal U)
\times\Gamma(\operatorname{Sym}^2T^*\mathcal U)
\label{eq:pb-two-channel-pair}
\end{equation}
for the sections \(q_i^b\) and \(q_i^m\), and an analogous pair of cubic tensors. The pair is
determined by the two statistical models, sections, and connections; a single scalar geometry is not
part of those data.

\definitionheading{Weighted product geometry}{hyp:pb-weighted-product-geometry}
If one joint geometry is required, positive constants \(w_b,w_m\) are declared and
\begin{equation}
h^{\mathrm{prod}}
=w_bh_{i,b}^{\omega_b}+w_mh_{i,m}^{\omega_m}.
\label{eq:pb-weighted-channel-pullback}
\end{equation}
The weights fix the relative normalization of the two channels; unit weights are a specialization,
not a consequence of the common principal bundle.

\propositionheading{Radical of the weighted product}{prop:pb-product-radical}
Under \Cref{hyp:pb-regular-models} in both channels,
\begin{equation}
\operatorname{rad}h^{\mathrm{prod}}
=\ker D^{\omega_b}q_i^b\cap\ker D^{\omega_m}q_i^m,
\label{eq:pb-product-radical}
\end{equation}
so two degenerate channel tensors are jointly nondegenerate exactly when their null distributions
intersect in zero. \status{ESTABLISHED}

\noindent\emph{Proof.}
\(h^{\mathrm{prod}}(X,X)=w_bg_b^F(D^{\omega_b}q_i^bX,D^{\omega_b}q_i^bX)
+w_mg_m^F(D^{\omega_m}q_i^mX,D^{\omega_m}q_i^mX)\) is a sum of nonnegative terms with positive
weights, which vanishes exactly when both vertical jets vanish; for a positive-semidefinite form the
diagonal null set is the radical. \hfill$\square$

The common group acts on the two channels separately and supplies no bilinear cross tensor on
\(V\mathcal E_b\times V\mathcal E_m\); such a term needs a separately declared invariant bilinear
form, and a comparison of channel points needs \(\Phi\) or \(\widetilde\Phi\) of
\Cref{def:geo-cross-morphisms}. Neither the relative frame nor equal channel dimensions supplies
either object.

\section{Rank, radicals, and quotient geometry}
\label{sec:pb-rank-quotient}

\theoremheading{Constant-rank quotient theorem}{thm:pb-pullback-rank-quotient}
Let \(s:\mathcal U\to E\) be a smooth section under \Cref{hyp:pb-regular-models}. Pointwise,
\begin{equation}
\operatorname{rad}h_s^\omega
=\ker D^\omega s,
\qquad
\operatorname{rank}h_s^\omega
=\operatorname{rank}D^\omega s.
\label{eq:pb-pullback-radical-rank}
\end{equation}
If \(D^\omega s\) has constant rank on \(\mathcal U\), then \(K_s^\omega=\ker D^\omega s\) is a smooth
subbundle, and
\begin{equation}
Q_s^\omega=T\mathcal U/K_s^\omega
\label{eq:pb-pullback-quotient-bundle}
\end{equation}
carries a positive-definite metric and a symmetric cubic tensor defined by
\begin{align}
\bar h_s^\omega([X],[Y])
&=h_s^\omega(X,Y),
\label{eq:pb-pullback-quotient-metric}\\
\bar c_s^\omega([X],[Y],[Z])
&=c_s^\omega(X,Y,Z),
\label{eq:pb-pullback-quotient-amari}
\end{align}
and the induced map \(Q_s^\omega\to\operatorname{im}D^\omega s\) is an isometric vector-bundle
isomorphism. \status{ESTABLISHED}

The proof is in \Cref{app:base-proofs}. The quotient is a vector bundle, not automatically the tangent
bundle of a quotient manifold, and the following records what the manifold reading requires.

\propositionheading{Requirements for a quotient manifold}{prop:pb-quotient-manifold-requirements}
Assume the constant-rank case of \Cref{thm:pb-pullback-rank-quotient}. (i) If there is a smooth
Hausdorff manifold \(\overline{\mathcal U}\) and a surjective submersion
\(q_K:\mathcal U\to\overline{\mathcal U}\) whose fibers are the null leaves and whose tangent kernel is
\(K_s^\omega\), then
\begin{equation}
D^\omega s([X,Y])=0
\quad
\text{for all }X,Y\in\Gamma(K_s^\omega);
\label{eq:pb-null-involutivity}
\end{equation}
involutivity does not supply such a quotient, even at constant rank. (ii) Tensors on
\(\overline{\mathcal U}\) whose pullbacks are \(h_s^\omega\) and \(c_s^\omega\) exist only if
\begin{equation}
\mathcal L_Zh_s^\omega=0,
\qquad
\mathcal L_Zc_s^\omega=0
\quad\text{for every }Z\in\Gamma(K_s^\omega),
\label{eq:pb-null-basicness}
\end{equation}
and this basicness can fail although a smooth Hausdorff quotient exists. (iii) For any
\(g^F\)-compatible connection \(\nabla\) on \(s^*VE\), with the covariant Lie derivative
\begin{equation}
(\mathcal L_Z^\nabla\alpha)(X)
=\nabla_Z(\alpha(X))-\alpha([Z,X])
\label{eq:pb-covariant-lie-derivative}
\end{equation}
of an \(s^*VE\)-valued one-form \(\alpha\), every \(Z\in\Gamma(K_s^\omega)\) satisfies
\begin{equation}
(\mathcal L_Zh_s^\omega)(X,Y)
=g^F\bigl((\mathcal L_Z^\nabla D^\omega s)(X),D^\omega sY\bigr)
+g^F\bigl(D^\omega sX,(\mathcal L_Z^\nabla D^\omega s)(Y)\bigr),
\label{eq:pb-basicness-criterion}
\end{equation}
and the right side does not depend on the choice of \(\nabla\). \status{ESTABLISHED}

The proof, with a dense-leaf witness for (i) and a non-basic witness for (ii), is in
\Cref{app:base-proofs}. Constant rank is itself needed, and integrability of the radical is a further
condition.

\propositionheading{A constant-rank nonintegrable radical}{prop:pb-contact-null-counterexample}
For the normal location family, the trivial principal translation bundle over \(\R^3\) with
coordinates \((x,y,z)\), the constant section \(\mu=0\), and the local connection form
\begin{equation}
A=\mathrm dz-x\mathrm dy,
\label{eq:pb-contact-connection}
\end{equation}
one has \(D^\omega s=A\),
\begin{equation}
h_s^\omega=(\mathrm dz-x\mathrm dy)^2,
\label{eq:pb-contact-pullback}
\end{equation}
of constant rank one, and
\begin{equation}
A\wedge\mathrm dA
=-\mathrm dz\wedge\mathrm dx\wedge\mathrm dy\ne0,
\label{eq:pb-contact-frobenius}
\end{equation}
so \(\ker D^\omega s\) is a contact distribution and not a foliation. Without constant rank the
quotient bundle fails: for the same family over \(\mathcal C=\R\) with zero connection and
\(s(x)=\mathcal N(x^2,1)\),
\begin{equation}
h_s=4x^2dx^2
\label{eq:pb-rank-jump-example}
\end{equation}
has rank one away from the origin and rank zero at it. \status{ESTABLISHED}

\noindent\emph{Proof.}
For a section of means \(\mu\), the vertical jet is \((d\mu+A)\partial_\mu\) by the sign convention of
\Cref{def:pb-connection-difference}, and the location Fisher metric is \(1\). Then
\(\mathrm dA=-\mathrm dx\wedge\mathrm dy\) and \(x\mathrm dy\wedge\mathrm dx\wedge\mathrm dy=0\), and
the Frobenius criterion identifies \(\ker A\) as nonintegrable. In the second example the vertical jet
is \(d\mu=2x\mathrm dx\), whose square is \eqref{eq:pb-rank-jump-example}. \hfill$\square$

\section{Transported divergence along base curves}
\label{sec:pb-transported-divergence}

Let \(\mathscr D\) be a contrast satisfying \Cref{hyp:pb-contrast}, with induced metric
\(g_{\mathscr D}\) of \eqref{eq:pb-divergence-metric-jet}, mixed third jets \(\Gamma_{\mathscr D}\)
and \(\Gamma^*_{\mathscr D}\) of \cref{eq:pb-divergence-first-connection-jet,eq:pb-divergence-dual-connection-jet},
and cubic tensor \(T^F_{\mathscr D}=\Gamma_{\mathscr D}-\Gamma^*_{\mathscr D}\) of
\eqref{eq:pb-divergence-amari-jet}. For a smooth base curve \(\gamma\) with \(\gamma(0)=c\), let
\(\Omega^\omega_{0\to\epsilon}\) be the channel transport of \Cref{def:geo-connections} along
\(\gamma|_{[0,\epsilon]}\), and bring the section value back to the initial fiber:
\begin{equation}
\widehat s_\gamma(\epsilon)
=(\Omega^\omega_{0\to\epsilon})^{-1}s(\gamma(\epsilon)).
\label{eq:pb-transported-section-curve}
\end{equation}

\corollaryheading{Quadratic transported divergence}{cor:pb-transported-divergence-quadratic}
The transported section curve has initial velocity
\begin{equation}
\left.\frac{d}{d\epsilon}\widehat s_\gamma(\epsilon)\right|_{\epsilon=0}
=D^\omega s(\dot\gamma(0)),
\label{eq:pb-transported-section-velocity}
\end{equation}
and with \(h_{s,\mathscr D}^\omega=(D^\omega s)^*g_{\mathscr D}\),
\begin{equation}
\mathscr D(s(c),\widehat s_\gamma(\epsilon))
=\frac{\epsilon^2}{2}h_{s,\mathscr D}^\omega
(\dot\gamma(0),\dot\gamma(0))
+O(\epsilon^3).
\label{eq:pb-transported-divergence-expansion}
\end{equation}
If \(g_{\mathscr D}=g^F\), as for the regular KL contrast of
\Cref{prop:pb-kl-divergence-jets}, the two pullbacks coincide,
\(h_{s,\mathscr D}^\omega=h_s^\omega\). \status{ESTABLISHED}

\noindent\emph{Proof.}
A horizontal curve over \(\gamma\) is carried by \((\Omega^\omega_{0\to\epsilon})^{-1}\) to a constant,
so the derivative of \(\widehat s_\gamma\) at zero is the vertical part of \(Ts(\dot\gamma(0))\),
which is \eqref{eq:pb-transported-section-velocity}. A contrast and its first derivative vanish on the
diagonal and its quadratic Taylor term in the second argument is the induced metric, as in the proof
of \Cref{cor:pb-orbit-fisher-form}; substitution gives the expansion. \hfill$\square$

The cubic coefficient of such a one-sided expansion is not determined by
\(c^\omega_{s,\mathscr D}:=(D^\omega s)^*T^F_{\mathscr D}\) alone, and it can be computed exactly. In a
chart write \(\Gamma_{\mathscr D,ij,k}=\Gamma_{\mathscr D}(\partial_i,\partial_j,\partial_k)\) and
\(\Gamma^*_{\mathscr D,ij,k}=\Gamma^*_{\mathscr D}(\partial_i,\partial_j,\partial_k)\), and let
\(\nabla^*_{\mathscr D}\) be the torsion-free connection on \(\mathcal B\) with
\(g_{\mathscr D}(\nabla^*_{\mathscr D,\partial_i}\partial_j,\partial_k)=\Gamma^*_{\mathscr D,ij,k}\).

\propositionheading{Cubic coefficient of a one-sided contrast expansion}{prop:pb-divergence-cubic-coefficient}
The connection \(\nabla^*_{\mathscr D}\) is invariant under the represented action and descends to
each fiber of \(E\). For every \(C^3\) curve \(c\) in \(\mathcal B\) with \(c(0)=p\),
\(\dot c(0)=v\), and covariant acceleration \(a^*=(\nabla^*_{\mathscr D,\dot c}\dot c)(0)\),
\begin{equation}
\mathscr D(p,c(\epsilon))
=\frac{\epsilon^2}{2}g_{\mathscr D}(v,v)
+\epsilon^3\left[\frac12g_{\mathscr D}(v,a^*)
+\frac16T^F_{\mathscr D}(v,v,v)\right]
+o(\epsilon^3).
\label{eq:pb-divergence-cubic-coefficient}
\end{equation}
For the transported section curve, \(v=D^\omega s(\dot\gamma(0))\) and the cubic coefficient in
\eqref{eq:pb-transported-divergence-expansion} is
\(\tfrac12g_{\mathscr D}(v,a^*)+\tfrac16c^\omega_{s,\mathscr D}(\dot\gamma(0),\dot\gamma(0),\dot\gamma(0))\)
with \(a^*\) the \(\nabla^*_{\mathscr D}\)-acceleration of \(\widehat s_\gamma\); it equals
\(\tfrac16c^\omega_{s,\mathscr D}\) exactly when \(a^*\) is \(g_{\mathscr D}\)-orthogonal to \(v\), for
instance along a \(\nabla^*_{\mathscr D}\)-geodesic. For regular KL,
\(\Gamma^*_{\mathscr D,ij,k}=\E_\theta[\mathbf s_k\partial_i\mathbf s_j]\) with
\(\mathbf s_i=\partial_i\log p_\theta\), the lowered coefficients of the exponential connection, which
vanish in the natural chart of an exponential family, and \(T^F_{\mathscr D}\) is the
Amari--Chentsov tensor. \status{ESTABLISHED}

\noindent\emph{Proof.}
Invariance of \(\mathscr D\) under the simultaneous represented action makes its mixed jets, hence
\(\nabla^*_{\mathscr D}\), invariant. Differentiating the diagonal identity
\(\partial'_i\partial'_j\mathscr D(\theta,\theta)=(g_{\mathscr D})_{ij}(\theta)\), obtained in the proof
of \Cref{cor:pb-orbit-fisher-form}, gives
\(\partial'_i\partial'_j\partial'_k\mathscr D=\partial_k(g_{\mathscr D})_{ij}+\Gamma^*_{\mathscr D,ij,k}\)
on the diagonal, and differentiating
\((g_{\mathscr D})_{ij}(\theta)=-\partial_i\partial'_j\mathscr D(\theta,\theta)\) gives
\(\partial_k(g_{\mathscr D})_{ij}=\Gamma_{\mathscr D,ki,j}+\Gamma^*_{\mathscr D,kj,i}\). The third
derivative of \(\epsilon\mapsto\mathscr D(p,c(\epsilon))\) at zero is
\((\partial'_i\partial'_j\partial'_k\mathscr D)v^iv^jv^k+3(g_{\mathscr D})_{ij}v^i\ddot c^j(0)\), which by
these identities equals
\(T^F_{\mathscr D}(v,v,v)+3\Gamma^*_{\mathscr D}(v,v,v)+3g_{\mathscr D}(v,\ddot c(0))
=T^F_{\mathscr D}(v,v,v)+3g_{\mathscr D}(v,a^*)\), because
\(a^{*k}=\ddot c^k(0)+\Gamma^{*k}_{\mathscr D,ij}v^iv^j\) with the index raised by
\(g_{\mathscr D}^{-1}\). The lower-order coefficients are those of
\Cref{cor:pb-transported-divergence-quadratic}, and Taylor's theorem gives
\eqref{eq:pb-divergence-cubic-coefficient}. For regular KL the last line of
\eqref{eq:pb-kl-jet-calculation} gives \(\Gamma^*_{\mathscr D}\) \citep{AmariNagaoka2000}.
\hfill$\square$

The divergence-induced cubic tensor is therefore the difference of the two mixed third jets, not an
arbitrary one-sided cubic Taylor coefficient.

\section{Morphisms over base maps}
\label{sec:pb-coarse-naturality}

Consider smooth associated bundles \(E\to\mathcal C\) and \(\bar E\to\bar{\mathcal C}\) with
connections \(\omega\) and \(\bar\omega\), a smooth map \(f:\mathcal C\to\bar{\mathcal C}\), and a
smooth bundle morphism \(\Psi:E\to\bar E\) over \(f\), with vertical differential
\(T^V\Psi:VE\to V\bar E\). For a section \(s\), the composite \(\Psi\circ s\) is only a section of the
pullback bundle \(f^*\bar E\).

\definitionheading{Related sections}{def:pb-related-sections}
Sections \(s\) of \(E\) and \(\bar s\) of \(\bar E\) are related by \(\Psi\) when
\begin{equation}
\Psi\circ s=\bar s\circ f.
\label{eq:pb-coarse-related-sections}
\end{equation}
For a surjective submersion \(f\), this excludes fine sections whose \(\Psi\)-images vary along a
fiber of \(f\).

\definitionheading{Horizontal defect of a bundle morphism}{def:pb-horizontal-defect}
Following the covariant defect \eqref{eq:geo-nonlinear-defect-phi}, the horizontal defect of \(\Psi\)
at \(e\in E_c\) is
\begin{equation}
(\mathcal D\Psi)_e(X)
=A_\Psi(e;X)
:=T_e\Psi(H_e^\omega X)
-H_{\Psi(e)}^{\bar\omega}(T_cfX),
\qquad X\in T_c\mathcal C,
\label{eq:pb-coarse-horizontal-defect}
\end{equation}
with the transported fine horizontal lift first and the coarse horizontal lift subtracted. Its type is
\(A_\Psi\in\Gamma(E;\varpi^*T^*\mathcal C\otimes\Psi^*V\bar E)\).

Both terms of \eqref{eq:pb-coarse-horizontal-defect} project to \(T_cfX\), so the defect is vertical
and equals \(\operatorname{ver}^{\bar\omega}(T_e\Psi(H_e^\omega X))\). Along a section,
\(A_\Psi(s;\cdot)\in\Gamma(\mathcal U;T^*\mathcal U\otimes f^*\bar s^*V\bar E)\).

\theoremheading{Sharp section descent}{thm:pb-section-descent}
Let \(f\) be a surjective smooth submersion, \(\Psi\) a smooth bundle morphism over \(f\), and
\(Q\in\Gamma(\mathcal C,E)\). Consider (P1) there is \(\bar Q\in\Gamma(\bar{\mathcal C},\bar E)\) with
\(\Psi\circ Q=\bar Q\circ f\); (P2) \(\Psi\circ Q\) is constant on each fiber of \(f\); and
\begin{equation}
\textbf{(P3)}\qquad
T^V\Psi\left(D^\omega Q(X)\right)+A_\Psi\left(Q(c);X\right)=0
\quad\text{for every }c\text{ and every }X\in\ker T_cf.
\label{eq:pb-projectability-criterion}
\end{equation}
Then (P1) and (P2) are equivalent, both imply (P3), and (P3) implies (P2) when the fibers of \(f\)
are connected. Under (P2) the descended \(\bar Q\) is unique, smooth, and a section. The submersion
hypothesis cannot be dropped: for \(f(x)=x^3\) on \(\R\), the trivial bundle with fiber
\(\{\mathcal N(\mu,1)\}\), \(\Psi\) the identity on fibers, and \(Q(x)=\mathcal N(x,1)\), the unique
descent \(\bar Q(y)=\mathcal N(y^{1/3},1)\) is not differentiable at \(y=0\). \status{ESTABLISHED}

\noindent\emph{Proof.}
That (P1) implies (P2) is immediate. A surjective smooth submersion is a smooth quotient map, so a
smooth map constant on its fibers descends uniquely to a smooth map; surjectivity gives uniqueness,
and \(\bar\varpi\bar Qf=\bar\varpi\Psi Q=f\varpi Q=f\) gives \(\bar\varpi\bar Q=\operatorname{id}\).
For (P2) implying (P3), take \(X\in\ker T_cf\), split \(T_cQX=H^\omega_{Q(c)}X+D^\omega Q(X)\), apply
\(T\Psi\), and use \(T_cfX=0\): what remains is
\(T(\Psi Q)(X)=A_\Psi(Q(c);X)+T^V\Psi(D^\omega Q(X))\), which fiber constancy forces to vanish. The
same computation shows that (P3) makes \(T(\Psi Q)\) annihilate \(\ker Tf\), and a smooth map with
vanishing derivative along a connected embedded submanifold is constant on it. In the example
\(f\) is a bijection that fails to be a submersion at \(0\). \hfill$\square$

\propositionheading{A pointwise bundle morphism is not a map on configurations}{prop:pb-nonfunctorial-descent}
Suppose \(f\) is a surjective submersion, \(\ker T_{c_0}f\ne0\) at some \(c_0\), there are
\(e_0\in E_{c_0}\) and \(w\in V_{e_0}E\) with \(T^V\Psi(w)\ne0\), and a global smooth section
\(Q_0\) satisfies \(Q_0(c_0)=e_0\). Then the projectable set \(\Gamma_{\mathrm{proj}}(\Psi)\) of
sections satisfying (P2) is a proper subset of \(\Gamma(\mathcal C,E)\). For the total collapse
\(S^1\to\{\ast\}\) with trivial bundles carrying \(\{\mathcal N(\mu,1)\}\), \(\Psi\) the identity on
fibers, and \(Q(x)=\mathcal N(\sin x,1)\), no coarse section exists, and in the linear realization on
\(L^2(S^1)\) the descendable set is the constants, a closed subspace of infinite codimension. If
\(f\) is a diffeomorphism every section descends. \status{ESTABLISHED}

\noindent\emph{Proof.}
Choose a chart in which \(\ker Tf\) is spanned by \(\partial_1,\ldots,\partial_k\) with \(k\geq1\). If
\(Q_0\) fails (P2) there is nothing to prove. Otherwise let \(W\) be a smooth vertical field in a local
trivialization with \(W(Q_0(c_0))=w\) and local flow \(\Phi^W\), take \(\chi\in C^\infty_c\) supported
in the chart with \(\chi(c_0)=0\) and \(\partial_1\chi(c_0)\ne0\), and set
\(Q_\epsilon(c)=\Phi^W_{\epsilon\chi(c)}(Q_0(c))\) on the support and \(Q_\epsilon=Q_0\) off it. Since
\(\chi(c_0)=0\), \(A_\Psi(Q_\epsilon(c_0);\partial_1)\) and \(T^V\Psi\) at that point do not depend on
\(\epsilon\), and differentiating the left side of \eqref{eq:pb-projectability-criterion} at
\(\epsilon=0\) gives \(\partial_1\chi(c_0)T^V\Psi(w)\ne0\), so (P3) fails for small
\(\epsilon\ne0\). In the collapse example \(A_\Psi=0\), \(\ker Tf=TS^1\), and
\(T^V\Psi(D^\omega Q\partial_x)=\cos(x)\partial_\mu\ne0\); a descendable section is constant on
\(S^1\). A diffeomorphism has trivial kernel and every fiber is a point. \hfill$\square$

\theoremheading{Covariant first-jet chain rule}{thm:pb-covariant-jet-naturality}
For sections related by \eqref{eq:pb-coarse-related-sections},
\begin{equation}
D^{\bar\omega}\bar s(TfX)
=T^V\Psi(D^\omega sX)+(\mathcal D\Psi)_{s(c)}(X).
\label{eq:pb-covariant-jet-chain-rule}
\end{equation}
Hence the square of covariant first jets commutes exactly when \(\mathcal D\Psi\) vanishes along
\(s\), and then every vertical covariant tensor \(\bar\tau\) on \(\bar E\) satisfies
\begin{equation}
f^*[(D^{\bar\omega}\bar s)^*\bar\tau]
=(D^\omega s)^*[(T^V\Psi)^*\bar\tau].
\label{eq:pb-covariant-tensor-naturality}
\end{equation}
\status{ESTABLISHED}

\noindent\emph{Proof.}
Differentiate \eqref{eq:pb-coarse-related-sections} and decompose
\(TsX=H_s^\omega X+D^\omega sX\). Applying \(T\Psi\) to the horizontal term gives the coarse horizontal
lift plus \((\mathcal D\Psi)_{s(c)}(X)\) by \eqref{eq:pb-coarse-horizontal-defect}; applying it to the
vertical term gives \(T^V\Psi(D^\omega sX)\). The vertical part of the other side is
\(D^{\bar\omega}\bar s(TfX)\). When the defect vanishes, applying \(\bar\tau\) to the intertwined jets
gives \eqref{eq:pb-covariant-tensor-naturality}. \hfill$\square$

\section{The base Fisher comparison}
\label{sec:pb-base-fisher-comparison}

A bundle morphism \(\Psi\) over \(f\) is \emph{fiberwise Markov} when \Cref{hyp:pb-regular-models}
holds at both scales and each fiber map \(E_c\to\bar E_{f(c)}\), read in frames at \(c\) and
\(f(c)\), satisfies \Cref{hyp:pb-markov-morphism}. The coarse metric is invariant under the
represented action (\Cref{prop:pb-statistical-tensor-descent}), so \((T^V\Psi)^*\bar g^F\) does not
depend on the frames, and the vertical Fisher defect \(\Delta_F^\Psi=g^F-(T^V\Psi)^*\bar g^F\) of
\eqref{eq:pb-vertical-fisher-defect} is defined along \(\Psi\) for every smooth \(\Psi\).

\definitionheading{Base comparison tensors}{def:pb-base-comparison-tensors}
Along related sections write \(\mathsf u_X=D^\omega sX\) for the fine vertical jet, \(\Psi_*=T^V\Psi\),
\(\mathsf a_X=A_\Psi(s(c);X)\) for the horizontal defect, and
\(\bar{\mathsf u}_X=D^{\bar\omega}\bar s(T_cfX)=\Psi_*\mathsf u_X+\mathsf a_X\) for the coarse vertical
jet, by \eqref{eq:pb-covariant-jet-chain-rule}. Define
\begin{align}
\delta_\Psi(X,Y)&=\Delta_F^\Psi(\mathsf u_X,\mathsf u_Y),
\label{eq:pb-base-fisher-defect}\\
\mathcal X_\Psi(X,Y)&=\bar g^F(\Psi_*\mathsf u_X,\mathsf a_Y)+\bar g^F(\mathsf a_X,\Psi_*\mathsf u_Y),
\label{eq:pb-cross-anomaly-tensor}\\
\mathcal Q_\Psi(X,Y)&=\bar g^F(\mathsf a_X,\mathsf a_Y).
\label{eq:pb-quadratic-anomaly-tensor}
\end{align}

\corollaryheading{Base Fisher contraction}{cor:pb-base-fisher-contraction}
Let \(\Psi\) be fiberwise Markov, let \(s,\bar s\) be related, and suppose \(\mathcal D\Psi=0\) along
\(s\). Then
\begin{equation}
h_s^\omega-f^*h_{\bar s}^{\bar\omega}
=(D^\omega s)^*\Delta_F^\Psi\succeq0,
\label{eq:pb-pullback-fisher-contraction}
\end{equation}
with equality on \(X\) exactly when the fine score in the direction \(D^\omega sX\) is measurable with
respect to the coarse output. \status{ESTABLISHED}

\noindent\emph{Proof.}
In a frame pair \(\Delta_F^\Psi\) is the vertical Fisher defect of the fiber map, so
\Cref{thm:pb-pullback-fisher-defect} gives \(\Delta_F^\Psi\succeq0\) with the stated equality case
\eqref{eq:pb-fisher-defect-score-variance}. With vanishing defect,
\eqref{eq:pb-covariant-tensor-naturality} for \(\bar\tau=\bar g^F\) gives
\(f^*h_{\bar s}^{\bar\omega}=(D^\omega s)^*(T^V\Psi)^*\bar g^F\), and subtracting this from
\eqref{eq:pb-fisher-pullback} proves \eqref{eq:pb-pullback-fisher-contraction}. \hfill$\square$

\theoremheading{Exact signed base Fisher comparison}{thm:pb-signed-base-comparison}
For sections related by \eqref{eq:pb-coarse-related-sections}, with no hypothesis on the defect,
\begin{equation}
h_s^\omega-f^*\bar h_{\bar s}^{\bar\omega}
=\delta_\Psi-\mathcal X_\Psi-\mathcal Q_\Psi
\label{eq:pb-signed-base-comparison}
\end{equation}
as symmetric two-tensors on the fine domain. Here \(\mathcal X_\Psi\) is symmetric and sign
indefinite and \(\mathcal Q_\Psi\succeq0\), and both enter with a minus sign. \status{ESTABLISHED}

\noindent\emph{Proof.}
By \eqref{eq:pb-covariant-jet-chain-rule},
\(f^*\bar h(X,Y)=\bar g^F(\Psi_*\mathsf u_X+\mathsf a_X,\Psi_*\mathsf u_Y+\mathsf a_Y)
=\bar g^F(\Psi_*\mathsf u_X,\Psi_*\mathsf u_Y)+\mathcal X_\Psi(X,Y)+\mathcal Q_\Psi(X,Y)\), and
\(\bar g^F(\Psi_*\mathsf u_X,\Psi_*\mathsf u_Y)=g^F(\mathsf u_X,\mathsf u_Y)-\delta_\Psi(X,Y)\).
Subtracting from \(h_s^\omega(X,Y)=g^F(\mathsf u_X,\mathsf u_Y)\) proves the identity.
\hfill$\square$

\theoremheading{Exact signed positivity criterion}{thm:pb-signed-positivity-criterion}
Let \(\Psi\) be fiberwise Markov, so that \(\delta_\Psi\succeq0\). At a fixed \(c\) the following are
equivalent: \(h_s^\omega-f^*\bar h_{\bar s}^{\bar\omega}\succeq0\); the coarse jet is no longer than the
fine jet, \(\lVert\bar{\mathsf u}_X\rVert_{\bar g^F}\leq\lVert\mathsf u_X\rVert_{g^F}\) for every
\(X\); and
\begin{equation}
2\bar g^F\left(\Psi_*\mathsf u_X,\mathsf a_X\right)
+\left\lVert\mathsf a_X\right\rVert^2_{\bar g^F}
\leq\delta_\Psi(X,X)
\quad\text{for every }X.
\label{eq:pb-positivity-criterion}
\end{equation}
Vanishing of the horizontal defect is therefore sufficient for base positivity and not necessary, and
the triangle-inequality margin
\(\lVert\mathsf a_X\rVert_{\bar g^F}\leq\sqrt{h(X,X)}-\sqrt{h(X,X)-\delta_\Psi(X,X)}\) is sufficient
and not necessary. Both failure directions occur: over \(\mathcal C=\bar{\mathcal C}=\R\) with
\(f=\operatorname{id}\), fine fiber \(\{\mathcal N(\mu,1)\}\), the channel adding independent
\(\mathcal N(0,1)\) noise, the section \(s(x)=\mathcal N(x,1)\), zero source connection, and a target
connection contributing \(\mathsf a_X=\mathrm c\partial_\mu\), the comparison equals
\(1-\tfrac12(1+\mathrm c)^2\), which is \(\tfrac{79}{200}\) at \(\mathrm c=\tfrac1{10}\) and
\(\tfrac{23}{25}\) at \(\mathrm c=-\tfrac35\), with nonzero defect in both cases and the margin violated
in the second, and \(-\tfrac18\) at \(\mathrm c=\tfrac12\); with the identity kernel, constant related
sections, source horizontal field \(\partial_x\), and target horizontal field
\(\partial_x+\mathrm c\partial_\mu\) with \(\mathrm c\ne0\), the comparison is
\(-\mathrm c^2dx^2\prec0\) at zero information loss. \status{ESTABLISHED}

\noindent\emph{Proof.}
The first two conditions are \eqref{eq:pb-signed-base-comparison} on the diagonal, with
\(h(X,X)=\lVert\mathsf u_X\rVert^2\) and \(f^*\bar h(X,X)=\lVert\Psi_*\mathsf u_X+\mathsf a_X\rVert^2\).
Expanding the square and substituting
\(\lVert\Psi_*\mathsf u_X\rVert^2=\lVert\mathsf u_X\rVert^2-\delta_\Psi(X,X)\) gives
\eqref{eq:pb-positivity-criterion}. In the first example \(g^F=1\), \(\bar g^F=\tfrac12\),
\(\Delta_F^\Psi=\tfrac12\), \(\mathsf u=1\), and \(\bar{\mathsf u}=1+\mathrm c\); at
\(\mathrm c=-\tfrac35\), \(\lVert\mathsf a\rVert=\tfrac{3}{5\sqrt2}\) exceeds \(1-\tfrac1{\sqrt2}\).
In the second, \(\mathsf u=0\) and \(\mathsf a=-\mathrm c\partial_\mu\). \hfill$\square$

Without the guard \(\mathcal D\Psi=0\), the vertical conclusion of
\Cref{thm:pb-pullback-fisher-defect} is unchanged, while
\eqref{eq:pb-pullback-fisher-contraction} is replaced by \eqref{eq:pb-signed-base-comparison} and its
criterion \eqref{eq:pb-positivity-criterion}.

When weighted products (\Cref{hyp:pb-weighted-product-geometry}) at two scales are compared, two
further declarations are consumed. Writing \(\mu_{\mathcal C}\) and \(\bar\mu_{\bar{\mathcal C}}\)
for finite positive base measures used to integrate the fine and coarse tensors and \(w_x,\bar w_x\)
for the two sets of channel weights, they are
\begin{equation}
\textbf{(X1)}\quad f_\#\mu_{\mathcal C}=\bar\mu_{\bar{\mathcal C}},
\qquad\qquad
\textbf{(X2)}\quad \bar w_x\circ f\leq w_x
\quad\text{for }x\in\{b,m\}.
\label{eq:pb-cross-scale-declarations}
\end{equation}

\propositionheading{Cross-scale comparison of weighted products}{prop:pb-weighted-cross-scale}
Under a channelwise comparison, with \(h_x\) and \(\bar h_x\) the channel pullbacks,
\begin{equation}
h^{\mathrm{prod}}-f^*\bar h^{\mathrm{prod}}
=\sum_{x\in\{b,m\}}
\left[
w_x\left(h_x-f^*\bar h_x\right)
+\left(w_x-\bar w_x\circ f\right)f^*\bar h_x
\right].
\label{eq:pb-two-channel-comparison}
\end{equation}
If each channel is fiberwise Markov with vanishing horizontal defect, so that
\(\delta_x:=h_x-f^*\bar h_x\succeq0\) by \Cref{cor:pb-base-fisher-contraction}, then (X2) makes the
difference positive semidefinite, and without (X2) positivity fails on \(X\) exactly when
\begin{equation}
\sum_{x\in\{b,m\}}(\bar w_x\circ f-w_x)f^*\bar h_x(X,X)
>\sum_{x\in\{b,m\}}w_x\delta_x(X,X);
\label{eq:pb-weight-positivity-failure}
\end{equation}
for one channel with \(f^*\bar h_x(X,X)=\delta_x(X,X)=\tfrac12\), as for the unit-variance normal
location section followed by additive unit-variance normal noise, unit fine weight gives
\(1-\tfrac12\bar w_x\), positive at \(\bar w_x=\tfrac32\) and negative at \(\bar w_x=3\). For
\(f\)-related vector fields \(X\) and \(\bar X\) with finite integrals, (X1) gives
\[
\int h^{\mathrm{prod}}(X,X)\mathrm d\mu_{\mathcal C}-\int\bar h^{\mathrm{prod}}(\bar X,\bar X)\mathrm d\bar\mu_{\bar{\mathcal C}}
=\int(h^{\mathrm{prod}}-f^*\bar h^{\mathrm{prod}})(X,X)\mathrm d\mu_{\mathcal C},
\]
which is nonnegative under the preceding hypotheses and (X2); without (X1) the left side can be
negative with vanishing fiberwise defect: for \(\mathcal C=\bar{\mathcal C}\) the disjoint union of two
copies of \(\R\), \(f=\operatorname{id}\), \(h\) zero on the first copy and \(dx^2\) on the second, and
unit point masses \(\mu_{\mathcal C}\) and \(\bar\mu_{\bar{\mathcal C}}\) on the first and second copy,
the fine integral is \(0\) and the coarse one \(1\). \status{ESTABLISHED}

\noindent\emph{Proof.}
Expand \(w_xh_x-(\bar w_x\circ f)f^*\bar h_x\). Under the hypotheses both summands in
\eqref{eq:pb-two-channel-comparison} are positive semidefinite when (X2) holds, and in general the
diagonal value is negative exactly under \eqref{eq:pb-weight-positivity-failure}. In the one-channel
witness the value is \(\tfrac12+(1-\bar w_x)\tfrac12\). For \(f\)-related fields
\(\bar h(\bar X,\bar X)\circ f=f^*\bar h(X,X)\), and the change-of-variables formula for
\(f_\#\mu_{\mathcal C}=\bar\mu_{\bar{\mathcal C}}\) moves the coarse integral to the fine domain.
\hfill$\square$

\corollaryheading{Coarse or meta-agent perceived geometry}{cor:pb-meta-perceived-geometry}
The coarse section \(\bar s\) carries its own connection-relative semimetric
\begin{equation}
\bar h_{\bar s}^{\bar\omega}
=(D^{\bar\omega}\bar s)^*\bar g^F.
\label{eq:pb-meta-perceived-geometry}
\end{equation}
For related sections with \(\mathcal D\Psi=0\) along \(s\),
\begin{equation}
f^*\bar h_{\bar s}^{\bar\omega}
=(D^\omega s)^*(T^V\Psi)^*\bar g^F,
\label{eq:pb-meta-perceived-naturality}
\end{equation}
and if \(\Psi\) is fiberwise Markov the fine-minus-coarse perceived-geometry defect is
\begin{equation}
h_s^\omega-f^*\bar h_{\bar s}^{\bar\omega}
=(D^\omega s)^*\Delta_F^\Psi\succeq0.
\label{eq:pb-meta-perceived-defect}
\end{equation}
\status{ESTABLISHED}

\noindent\emph{Proof.}
The first display is \Cref{def:pb-informational-pullbacks} at the coarse scale, the second is
\eqref{eq:pb-covariant-tensor-naturality} with \(\bar\tau=\bar g^F\), and the third is
\Cref{cor:pb-base-fisher-contraction}. \hfill$\square$

The comparison is contravariant: the coarse tensor is pulled back to the fine tangent space, and no
pushforward of a Riemannian metric is used. If the coarse operation is a restriction along an
embedding, an energy precomposition, or a parameter-dependent kernel, the Markov hypothesis fails and
the corollary does not apply. A Galerkin aggregation of a Gaussian law is the pushforward through a
precision-dependent map, which is a fixed kernel on each fixed-precision fiber and on no larger domain
(\Cref{prop:ig-galerkin-pushforward}); the corollary applies along such a fiber and not across
precisions.

When levelwise agent or meta-agent sections and regular vertical Fisher tensors are declared at every
level of a scale diagram, the geometric state may include their connection-relative perceived
semimetrics, which the associated bundles alone do not supply. A meta-agent thus has its own pullback
geometry, and its cross-scale naturality is \Cref{cor:pb-meta-perceived-geometry}. No tensor on one
level base is subtracted from a tensor on another before both are transported to a common reference
space.

\definitionheading{Perceived-geometry beta}{def:pb-metric-beta}
If the reference-space declarations of \Cref{ch:renormalization} include base diffeomorphisms
\(i_\ell:\mathcal C_\ell\to\mathcal C_\star\) and a block factor \(b_\ell>1\), the discrete
perceived-geometry beta on \(\mathcal C_\star\) is
\begin{equation}
\widetilde h_\ell=(i_\ell^{-1})^*h_\ell,
\qquad
\mathfrak B_{h,\ell}
=\frac{\widetilde h_{\ell+1}-\widetilde h_\ell}{\log b_\ell},
\label{eq:pb-metric-beta-reference}
\end{equation}
and on a differentiable reference-space curve \(t\mapsto\widetilde h_t\) its continuous counterpart
is \(\partial_t\widetilde h_t\). Both are conditional on the declared
identifications and scale normalization; without them \(h_{\ell+1}-h_\ell\) is not a well-typed
difference.

\section{The scale-connection defect and the isotropy criterion}
\label{sec:pb-isotropy}

The horizontal defect of an induced morphism is not an arbitrary vertical field. Suppose \(\Psi\) is
induced by an equivariant principal scale map \(\mathcal P:\mathscr P\to\bar{\mathscr P}\) over
\(f:\mathcal C\to\bar{\mathcal C}\) with Lie-group homomorphism \(\kappa:G\to\bar G\) and law-fiber map
\(q\), in the sense of \eqref{eq:rg-associated-scale-map} and
\Cref{prop:rg-principal-scale-map-existence}.

\definitionheading{Scale-connection defect form}{def:pb-scale-connection-defect}
The scale-connection defect form of \(\mathcal P\) is
\begin{equation}
\mathfrak A_{\mathcal P}
=\mathcal P^*\bar\omega-\mathrm d\kappa\circ\omega
\in\Omega^1(\mathscr P,\bar{\mathfrak g}).
\label{eq:pb-scale-connection-defect}
\end{equation}

\theoremheading{Isotropy criterion for a vanishing horizontal defect}{thm:pb-isotropy-criterion}
The form \(\mathfrak A_{\mathcal P}\) is horizontal and \(\operatorname{Ad}\circ\kappa\)-equivariant,
hence descends to an element of \(\Omega^1(\mathcal C;f^*\operatorname{Ad}\bar{\mathscr P})\), and
\begin{equation}
A_\Psi(e;X)
=\mathsf J_{\Psi(e)}\left(\mathfrak A_{\mathcal P}(X)\right),
\label{eq:pb-anomaly-fundamental-field}
\end{equation}
so the horizontal defect is the fundamental vertical field of a \(\bar{\mathfrak g}\)-valued one-form
on the base. For related sections, \(A_\Psi(s(c);X)=0\) for every \(c\) and every
\(X\in T_c\mathcal C\) if and only if
\begin{equation}
\mathfrak A_{\mathcal P}(X)\in\bar{\mathfrak g}_{\bar s(f(c))}
\quad\text{for every }c\text{ and every }X,
\label{eq:pb-isotropy-criterion}
\end{equation}
the isotropy subalgebra of the coarse section value. The principal identity
\(\mathfrak A_{\mathcal P}=0\) is sufficient; where the isotropy subalgebra is trivial, meaning the
action is locally free at \(\bar s(f(c))\), the criterion forces \(\mathfrak A_{\mathcal P}\) to vanish on
\(T_c\mathcal C\), and where it is nontrivial the criterion permits nonzero values in it. Moreover
\(\mathcal Q_\Psi(X,Y)=\mathfrak k_{\bar s(f(c))}(\mathfrak A_{\mathcal P}(X),\mathfrak A_{\mathcal P}(Y))\),
where \(\mathfrak k_{\bar e}(\xi,\eta)=\bar g^F(\mathsf J_{\bar e}\xi,\mathsf J_{\bar e}\eta)\) is the
orbit Fisher form \eqref{eq:pb-orbit-fisher-form} of the coarse fiber, positive semidefinite with
radical the isotropy subalgebra (\Cref{cor:pb-orbit-fisher-form}). \status{ESTABLISHED}

\noindent\emph{Proof.}
Evaluating \(\mathcal P^*\bar\omega\) on the fundamental field of \(\xi\) gives
\(\bar\omega(\bar{\mathsf f}_{\mathrm d\kappa\xi})=\mathrm d\kappa\xi\), so the difference annihilates vertical
vectors; equivariance follows from
\(R_g^*\mathcal P^*\bar\omega=\operatorname{Ad}_{\kappa(g)^{-1}}\mathcal P^*\bar\omega\) and
\(\mathrm d\kappa\circ\operatorname{Ad}_{g^{-1}}=\operatorname{Ad}_{\kappa(g)^{-1}}\circ\mathrm d\kappa\).
In local trivializations write \(\Psi(c,\beta)=(f(c),\psi_c\beta)\) with
\(\psi_c=\widehat{\bar\rho}(\varsigma(c))\circ q\). Using \(H^\omega_{(c,\beta)}X=(X,-\mathsf f_{A(X)}\beta)\)
and the differentiated intertwining \(Tq\circ\mathsf f_\xi=\bar{\mathsf f}_{\mathrm d\kappa(\xi)}\circ q\), which is
\eqref{eq:rg-scale-intertwiner} at \(g=\exp(t\xi)\), gives
\(A_\Psi(X)=\bar{\mathsf f}_{\mathfrak a(X)}(\psi_c\beta)\) with
\(\mathfrak a=\operatorname{Ad}_\varsigma(u^*\mathfrak A_{\mathcal P})\), which is
\eqref{eq:pb-anomaly-fundamental-field}. A fundamental field vanishes at \(\bar\beta\) exactly when its
generator lies in the isotropy subalgebra, and conjugation by \(\operatorname{Ad}_\varsigma\) moves
both together, which gives \eqref{eq:pb-isotropy-criterion}. The last statement is
\eqref{eq:pb-anomaly-fundamental-field} substituted into \eqref{eq:pb-quadratic-anomaly-tensor}.
\hfill$\square$

\definitionheading{Connection compatibility}{def:pb-connection-compatibility}
A principal scale map is compatible with the connections along related sections when
\eqref{eq:pb-isotropy-criterion} holds. No weaker or informal reading of compatibility is used in this
supplement.

\theoremheading{Ordered composition of horizontal defects}{thm:pb-anomaly-composition}
For composable morphisms \(E_0\xrightarrow{\Psi_{01}}E_1\xrightarrow{\Psi_{12}}E_2\) over
\(f_{01},f_{12}\), with \(\Psi_{02}=\Psi_{12}\circ\Psi_{01}\) and \(f_{02}=f_{12}\circ f_{01}\),
\begin{equation}
A_{\Psi_{02}}(e;X)
=T^V\Psi_{12}\big|_{\Psi_{01}(e)}
\left(A_{\Psi_{01}}(e;X)\right)
+A_{\Psi_{12}}\left(\Psi_{01}(e);T_cf_{01}X\right),
\label{eq:pb-anomaly-composition}
\end{equation}
and for induced morphisms, at the connection level,
\begin{equation}
\mathfrak A_{\mathcal P_{02}}
=\mathcal P_{01}^*\mathfrak A_{\mathcal P_{12}}
+\mathrm d\kappa_{12}\circ\mathfrak A_{\mathcal P_{01}}.
\label{eq:rg-scale-connection-defect-cocycle}
\end{equation}
Consequently \(A_{\Psi_{02}}(e;X)=0\) holds if and only if
\begin{equation}
T^V\Psi_{12}\big|_{\Psi_{01}(e)}A_{\Psi_{01}}(e;X)
=-A_{\Psi_{12}}\left(\Psi_{01}(e);T_cf_{01}X\right).
\label{eq:pb-anomaly-composition-vanishing}
\end{equation}
Factorwise vanishing is sufficient and not necessary: on the abelian three-level instance with
\(f_{01}(x)=2x\), \(f_{12}(y)=3y\), local connection forms \(0\), \(\mathrm c_1dy\), \(\mathrm c_2dz\), and
identity fiber maps, the composite defect is \(2\mathrm c_1+2(3\mathrm c_2-\mathrm c_1)=6\mathrm c_2\),
so \(\mathrm c_2=0\) and \(\mathrm c_1\ne0\) give a vanishing composite defect with both factor defects
nonzero. \status{ESTABLISHED}

\noindent\emph{Proof.}
Expand \(T\Psi_{01}(H^{\omega_0}_eX)=H^{\omega_1}_{\Psi_{01}e}(Tf_{01}X)+A_{\Psi_{01}}(e;X)\) and apply
\(T\Psi_{12}\). The horizontal summand contributes
\(H^{\omega_2}_{\Psi_{02}e}(Tf_{02}X)+A_{\Psi_{12}}(\Psi_{01}e;Tf_{01}X)\) by
\eqref{eq:pb-coarse-horizontal-defect}, and the vertical summand contributes
\(T^V\Psi_{12}(A_{\Psi_{01}}(e;X))\) because a bundle morphism preserves verticality; subtracting
\(H^{\omega_2}(Tf_{02}X)\) gives \eqref{eq:pb-anomaly-composition}. The connection-level form follows
from \(T^V\Psi_{12}\circ\mathsf J=\mathsf J\circ\mathrm d\kappa_{12}\), or directly from
\((\mathcal P_{12}\mathcal P_{01})^*\omega_2-\mathrm d\kappa_{12}\mathrm d\kappa_{01}\omega_0
=\mathcal P_{01}^*(\mathcal P_{12}^*\omega_2-\mathrm d\kappa_{12}\omega_1)
+\mathrm d\kappa_{12}(\mathcal P_{01}^*\omega_1-\mathrm d\kappa_{01}\omega_0)\). In the instance, the
first defect is \(2\mathrm c_1\mathrm dx\) and the second is \((3\mathrm c_2-\mathrm c_1)dy\), evaluated on the
pushed vector \(2\partial_y\). \hfill$\square$

The order in \eqref{eq:pb-anomaly-composition} is not cosmetic. The earlier defect takes values in
\(VE_1\) and the later one in \(VE_2\), so the unweighted sum of the two is a type error before it is
an error of fact: the earlier defect must be pushed by the later vertical differential, and the later
one evaluated at the image point on the pushed base vector. In the instance, dropping the base
pushforward gives \(\mathrm c_1+3\mathrm c_2\).

\begin{figure}[htbp]
\centering
\begin{tikzpicture}[>=stealth]
  \node[draw,rounded corners,fill=lightblue,minimum width=2.7cm,minimum height=0.9cm]
    (tc) at (0,2.1) {$T_c\mathcal C$};
  \node[draw,rounded corners,fill=lightgreen,minimum width=3.3cm,minimum height=0.9cm]
    (ve) at (5.4,2.1) {$V_{s(c)}E$};
  \node[draw,rounded corners,fill=lightblue,minimum width=2.7cm,minimum height=0.9cm]
    (tcb) at (0,0) {$T_{f(c)}\bar{\mathcal C}$};
  \node[draw,rounded corners,fill=lightgreen,minimum width=3.3cm,minimum height=0.9cm]
    (veb) at (5.4,0) {$V_{\bar s(f(c))}\bar E$};
  \draw[->,very thick,primaryblue] (tc)--node[above] {$D^\omega s$} (ve);
  \draw[->,very thick,primaryblue] (tcb)--node[below] {$D^{\bar\omega}\bar s$} (veb);
  \draw[->,very thick,sciencegreen] (tc)--node[left] {$T_cf$} (tcb);
  \draw[->,very thick,sciencegreen] (ve)--node[right] {$T^V\Psi$} (veb);
  \node[align=center,text width=7.0cm] at (2.7,-1.05)
    {$\mathcal D\Psi=0$: the first-jet square commutes};
\end{tikzpicture}
\caption{Naturality of the connection-split first jet. The square commutes exactly under the isotropy
criterion \eqref{eq:pb-isotropy-criterion}. When the horizontal defect is nonzero,
\eqref{eq:pb-covariant-jet-chain-rule} records the missing vertical term and
\eqref{eq:pb-signed-base-comparison} records the two retained base tensors \(-\mathcal X_\Psi\) and
\(-\mathcal Q_\Psi\).}
\label{fig:pb-pullback-naturality}
\end{figure}

\corollaryheading{Null directions under a parallel coarse morphism}{cor:pb-coarse-null-map}
If \(\mathcal D\Psi=0\) along \(s\), then
\begin{equation}
Tf(\ker D^\omega s)\subseteq\ker D^{\bar\omega}\bar s,
\label{eq:pb-coarse-null-inclusion}
\end{equation}
so a coarse morphism can create null directions, and first-jet naturality does not imply rank
preservation. \status{ESTABLISHED}

\noindent\emph{Proof.}
Set \(D^\omega sX=0\) in \eqref{eq:pb-covariant-jet-chain-rule} and use \(\mathcal D\Psi=0\).
\hfill$\square$

\section{The base Fisher defect cocycle}
\label{sec:pb-base-defect-cocycle}

For composable bundle morphisms the vertical cocycle \eqref{eq:pb-fisher-defect-cocycle} holds at every
point, since its proof is the contravariance of tensor pullback. Pulled back to the base it is not
additive in general, and the residual is available in closed form. Fix sections \(s_0,s_1,s_2\)
related at each stage, write \(\delta_{jk}=(D^{\omega_j}s_j)^*\Delta_F^{\Psi_{jk}}\) for the three base
defects, and along \(s_0\) write \(\mathsf p_X=T^V\Psi_{01}(D^{\omega_0}s_0X)\),
\(\mathsf a_X=A_{\Psi_{01}}(s_0;X)\), and
\(\bar{\mathsf u}_X=D^{\omega_1}s_1(Tf_{01}X)=\mathsf p_X+\mathsf a_X\).

\theoremheading{Base Fisher defect cocycle: exact residual and sharp criterion}{thm:pb-base-defect-cocycle}
The residual \(\mathcal R:=\delta_{02}-\delta_{01}-f_{01}^*\delta_{12}\) has the three equivalent exact
forms
\begin{align}
\mathcal R(X,Y)
&=\Delta_F^{\Psi_{12}}(\mathsf p_X,\mathsf p_Y)-\Delta_F^{\Psi_{12}}(\bar{\mathsf u}_X,\bar{\mathsf u}_Y),
\label{eq:pb-base-cocycle-residual-a}\\
&=-\left[\Delta_F^{\Psi_{12}}(\bar{\mathsf u}_X,\mathsf a_Y)
+\Delta_F^{\Psi_{12}}(\mathsf a_X,\bar{\mathsf u}_Y)\right]
+\Delta_F^{\Psi_{12}}(\mathsf a_X,\mathsf a_Y),
\label{eq:pb-base-cocycle-residual-b}\\
&=-\left[\Delta_F^{\Psi_{12}}(\mathsf p_X,\mathsf a_Y)
+\Delta_F^{\Psi_{12}}(\mathsf a_X,\mathsf p_Y)\right]
-\Delta_F^{\Psi_{12}}(\mathsf a_X,\mathsf a_Y).
\label{eq:pb-base-cocycle-residual-c}
\end{align}
The sign of the quadratic term is coupled to the argument of the cross terms, and mixing the two
conventions inflates the residual by exactly \(2\Delta_F^{\Psi_{12}}(\mathsf a_X,\mathsf a_Y)\). The
sharp base cocycle \(\delta_{02}=\delta_{01}+f_{01}^*\delta_{12}\) holds if and only if
\begin{equation}
\Delta_F^{\Psi_{12}}(\mathsf p_X,\mathsf p_X)
=\Delta_F^{\Psi_{12}}(\bar{\mathsf u}_X,\bar{\mathsf u}_X)
\quad\text{for every }X,
\label{eq:pb-base-cocycle-criterion}
\end{equation}
equivalently \(\Delta_F^{\Psi_{12}}(\mathsf a_X,2\mathsf p_X+\mathsf a_X)=0\) for every \(X\); if the
second arrow is fiberwise Markov this is equality of two \(\Delta_F^{\Psi_{12}}\)-seminorms.
\status{ESTABLISHED}

\noindent\emph{Proof.}
Applying \((D^{\omega_0}s_0)^*\) to \eqref{eq:pb-fisher-defect-cocycle} gives
\(\delta_{02}=\delta_{01}+\Delta_F^{\Psi_{12}}(\mathsf p_X,\mathsf p_Y)\), and by definition
\(f_{01}^*\delta_{12}(X,Y)=\Delta_F^{\Psi_{12}}(\bar{\mathsf u}_X,\bar{\mathsf u}_Y)\), which gives
\eqref{eq:pb-base-cocycle-residual-a}. Substituting \(\bar{\mathsf u}=\mathsf p+\mathsf a\) in the second
term gives \eqref{eq:pb-base-cocycle-residual-c}, and \(\mathsf p=\bar{\mathsf u}-\mathsf a\) in the
first gives \eqref{eq:pb-base-cocycle-residual-b}. A symmetric form vanishes exactly when it vanishes on
the diagonal. \hfill$\square$

\propositionheading{Sufficient conditions for the sharp base cocycle}{prop:pb-base-cocycle-sufficient}
Each of the following implies \eqref{eq:pb-base-cocycle-criterion}: the stage-one horizontal defect
vanishes, which by \Cref{thm:pb-isotropy-criterion} is \eqref{eq:pb-isotropy-criterion} at the first
arrow; the second arrow is Fisher lossless, \(\Delta_F^{\Psi_{12}}=0\); or
\(\mathsf a_X\in\operatorname{rad}\Delta_F^{\Psi_{12}}\) for every \(X\). The first two each imply the
third and are not ordered by strength, and none of the three is necessary. On the three-level
instance with normal location fibers of variance \(1,2,3\), sections \(x\mapsto x\), \(y\mapsto y/2\),
\(z\mapsto z/6\) of the means, base maps \(f_{01}(x)=2x\) and \(f_{12}(y)=3y\), additive normal noise
channels, and local connection forms \(0\), \(\mathrm c_1dy\), \(\mathrm c_2dz\), one has
\(\Delta_F^{\Psi_{12}}=\tfrac16\), stage-one defect \(2\mathrm c_1\), and
\(\mathcal R=-\tfrac23\mathrm c_1(\mathrm c_1+1)\), which vanishes at \(\mathrm c_1=-1\) where the
stage-one defect is \(-2\); the convention-mixed expression is \(\tfrac23\mathrm c_1(\mathrm c_1-1)\), and
at \(\mathrm c_1=\tfrac1{10}\) the two values are \(-\tfrac{11}{150}\) and \(-\tfrac3{50}\).
\status{ESTABLISHED}

\noindent\emph{Proof.}
Every term of \eqref{eq:pb-base-cocycle-residual-c} contains \(\Delta_F^{\Psi_{12}}\mathsf a\), and each
of the first two conditions makes it zero. At \(\mathrm c_1=0\) of the instance the stage-one defect
vanishes while \(\Delta_F^{\Psi_{12}}=\tfrac12-\tfrac13=\tfrac16\), and replacing its second arrow by the
identity gives \(\Delta_F^{\Psi_{12}}=0\) with stage-one defect \(2\mathrm c_1\ne0\) for
\(\mathrm c_1\ne0\). With the sign convention of \Cref{def:pb-connection-difference}, the jets are
\(\mathsf p=1\) and \(\bar{\mathsf u}=2(\tfrac12+\mathrm c_1)=1+2\mathrm c_1\) per \(dx\), so
\(\mathcal R=\tfrac16(1-(1+2\mathrm c_1)^2)\), and the mixed expression
\(-2\Delta_F^{\Psi_{12}}(\mathsf p,\mathsf a)+\Delta_F^{\Psi_{12}}(\mathsf a,\mathsf a)\) equals
\(\tfrac16(4\mathrm c_1^2-4\mathrm c_1)\). \hfill$\square$

The sharp base cocycle requires only the stage-one condition
\eqref{eq:pb-base-cocycle-criterion}. Separately, for each arrow \(jk\), the base defect
\(\delta_{jk}\) equals the information loss \(h_j-f_{jk}^*h_k\) exactly when
\(\mathcal X_{\Psi_{jk}}+\mathcal Q_{\Psi_{jk}}=0\), by
\Cref{thm:pb-signed-base-comparison}. Vanishing of the horizontal defect is sufficient for this
identification, not necessary: the signed cross term can cancel the quadratic term. When these
identifications hold, the sharp base cocycle is an additive decomposition of information loss on
the base.

\section{Intrinsic classification of connection defects}
\label{sec:pb-intrinsic-defect-classification}

The positivity criterion \eqref{eq:pb-positivity-criterion} and the cocycle criterion
\eqref{eq:pb-base-cocycle-criterion} are conditions on the horizontal defect, and by
\Cref{thm:pb-isotropy-criterion} the defect is the fundamental field of the connection defect. One can
therefore ask for all target connections that satisfy them. Fix the sections, the principal scale map,
the fiber morphism, and the source connection. At one base point write \(\mathcal V=T_c\mathcal C\),
\(E_0=V_{s(c)}E\), and \(E_1=V_{\bar s(f(c))}\bar E\) with Fisher inner products \(g_0\) and \(g_1\);
\(\mathsf D=D^\omega s\) at \(c\), \(\Psi_*=T^V\Psi\), and the pushed jet \(\mathsf p=\Psi_*\mathsf D\);
\(\mathfrak a_c=(f^*\operatorname{Ad}\bar{\mathscr P})_c\) and the fundamental-action map
\(\mathsf J=\mathsf J_{\bar s(f(c))}:\mathfrak a_c\to E_1\). A descended connection defect
\(\eta:\mathcal V\to\mathfrak a_c\) produces the anomaly \(\mathsf J\eta\), with image in the orbit
tangent \(\mathcal O=\operatorname{im}\mathsf J\).

\definitionheading{Relaxed pointwise defect problems}{def:pb-intrinsic-defect-problems}
The positivity problem and the cocycle problem are
\begin{align}
(\mathsf p+\mathsf J\eta)^*g_1(\mathsf p+\mathsf J\eta)
&\preceq\mathsf D^*g_0\mathsf D,
\label{eq:pb-intrinsic-positivity-problem}\\
(\mathsf p+\mathsf J\eta)^*\Delta(\mathsf p+\mathsf J\eta)
&=\mathsf p^*\Delta\mathsf p,
\label{eq:pb-intrinsic-cocycle-problem}
\end{align}
where \(\Delta\) in the second is the second-arrow vertical Fisher defect. They are separate
predicates; imposing both intersects their solution sets for the same \(\eta\). Varying the source
connection changes \(\mathsf D\) and is a different problem.

\propositionheading{Relaxation and the kernel constraint}{prop:pb-intrinsic-kernel-constraint}
Every defect realized by a target connection under related sections satisfies
\begin{equation}
(\mathsf p+\mathsf J\eta)|_{\ker T_cf}=0,
\label{eq:pb-intrinsic-kernel-constraint}
\end{equation}
which is condition (P3) of \Cref{thm:pb-section-descent} at \(c\). The problems of
\Cref{def:pb-intrinsic-defect-problems} do not impose it, so the parameterizations below classify the
relaxed pointwise problems, and a solution violating \eqref{eq:pb-intrinsic-kernel-constraint} is
realized by no lift. \status{ESTABLISHED}

\noindent\emph{Proof.}
By \Cref{thm:pb-covariant-jet-naturality}, \(D^{\bar\omega}\bar s(T_cfX)=(\mathsf p+\mathsf J\eta)X\),
and the left side vanishes for \(X\in\ker T_cf\). The reference defect \(\eta_0\) of
\Cref{thm:pb-defect-connection-descent} satisfies the constraint, so for a violating \(\eta\) the
difference \(\mathsf J(\eta-\eta_0)\) does not vanish on \(\ker T_cf\), and neither does
\(\eta-\eta_0\), which the annihilation clause of that theorem excludes for every lift.
\hfill$\square$

\theoremheading{Complete positivity parameterization}{thm:pb-defect-contractions}
Let \(P_{\mathcal O}\) be the \(g_1\)-orthogonal projection onto \(\mathcal O\) and set
\begin{equation}
\mathsf p_\perp=(I-P_{\mathcal O})\mathsf p,\qquad \mathsf p_{\mathcal O}=P_{\mathcal O}\mathsf p,\qquad
\mathsf H=\mathsf D^*g_0\mathsf D-\mathsf p_\perp^*g_1\mathsf p_\perp.
\label{eq:pb-defect-positivity-budget}
\end{equation}
The positivity problem \eqref{eq:pb-intrinsic-positivity-problem} has a pointwise solution if and only
if \(\mathsf H\succeq0\). In that case let
\(\operatorname{pr}_{\mathsf H}:\mathcal V\to\mathcal V_{\mathsf H}=\mathcal V/\operatorname{rad}\mathsf H\)
be the quotient with the inner product induced by \(\mathsf H\). All solutions are exactly
\begin{equation}
\mathsf J\eta=\mathsf k\operatorname{pr}_{\mathsf H}-\mathsf p_{\mathcal O},\qquad
\mathsf k:(\mathcal V_{\mathsf H},\mathsf H)\longrightarrow(\mathcal O,g_1)
\text{ a linear contraction},
\label{eq:pb-defect-contraction-parameterization}
\end{equation}
every such right side has a lift through \(\mathsf J\), and the lift ambiguity is
\(\operatorname{Hom}(\mathcal V,\ker\mathsf J)\). Zero-dimensional spaces and noninjective contractions
are included. \status{ESTABLISHED}

\corollaryheading{Markov first arrow: a convex feasible set}{cor:pb-defect-markov-feasible}
If the first arrow is fiberwise Markov, then
\(\delta=\mathsf D^*(g_0-\Psi_*^*g_1\Psi_*)\mathsf D\succeq0\) and
\begin{equation}
\mathsf H=\delta+\mathsf p_{\mathcal O}^*g_1\mathsf p_{\mathcal O}\succeq0,
\label{eq:pb-defect-markov-budget}
\end{equation}
so the feasible set of \eqref{eq:pb-intrinsic-positivity-problem} is nonempty, contains the zero
anomaly, and is convex, as is its set of lifts. These facts do not establish descent to a target
connection. \status{ESTABLISHED}

\theoremheading{Complete positive-semidefinite cocycle parameterization}{thm:pb-defect-partial-isometries}
Suppose \(\Delta\succeq0\). Let
\(\operatorname{pr}_\Delta:E_1\to E_\Delta=E_1/\operatorname{rad}\Delta\) be the metric quotient and define
\begin{equation}
\mathsf J_\Delta=\operatorname{pr}_\Delta\mathsf J,\qquad
\mathcal O_\Delta=\operatorname{im}\mathsf J_\Delta,\qquad
\mathsf p_\Delta=\operatorname{pr}_\Delta\mathsf p,\qquad
\mathsf w_\Delta=P_\Delta\mathsf p_\Delta,
\label{eq:pb-defect-cocycle-quotient-data}
\end{equation}
with \(P_\Delta\) the orthogonal projection onto \(\mathcal O_\Delta\). The solutions of
\eqref{eq:pb-intrinsic-cocycle-problem} are exactly
\begin{equation}
\mathsf J_\Delta\eta=(\mathsf T-\iota)\mathsf w_\Delta,\qquad
\mathsf T:\operatorname{im}\mathsf w_\Delta\longrightarrow\mathcal O_\Delta
\text{ a linear isometric embedding},
\label{eq:pb-defect-isometry-parameterization}
\end{equation}
with \(\iota\) the inclusion and lift ambiguity \(\operatorname{Hom}(\mathcal V,\ker\mathsf J_\Delta)\),
which contains both action isotropy and action directions invisible to \(\Delta\). The radical and
isotropy sufficient conditions correspond to \(\mathsf T=\iota\); \(\mathsf T=-\iota\) gives
\(\mathsf J_\Delta\eta=-2\mathsf w_\Delta\) without a vanishing anomaly or a lossless second arrow; if
\(\mathsf w_\Delta=0\) the equation forces \(\mathsf J_\Delta\eta=0\), whereas \(\Delta=0\) makes every
\(\eta\) a solution. Solutions realized under related sections additionally satisfy
\eqref{eq:pb-intrinsic-kernel-constraint}. \status{ESTABLISHED}

\theoremheading{Complete signed cocycle parameterization}{thm:pb-defect-signed-gram}
Let \(\Delta\) be any symmetric form on \(E_1\). Put
\(\mathcal O_0=\operatorname{rad}(\Delta|_{\mathcal O})\), choose a complement
\(\mathcal O=\mathcal O_1\oplus\mathcal O_0\), and write \(\Delta_1\) for the nondegenerate
restriction to \(\mathcal O_1\), with positive and negative indices \((n_+,n_-)\). Define
\(\mathsf w:\mathcal V\to\mathcal O_1\) and \(\mathsf l_x\in\mathcal O_0^*\) by
\begin{equation}
\Delta_1(\mathsf wx,o)=\Delta(\mathsf px,o)\quad(o\in\mathcal O_1),\qquad
\mathsf l_x(n)=\Delta(\mathsf px,n)\quad(n\in\mathcal O_0).
\label{eq:pb-defect-signed-projection}
\end{equation}
For a linear map \(\mathsf n:\mathcal V\to\mathcal O_0\) set
\begin{equation}
\mathsf G_{\mathsf n}(x,y)=\Delta_1(\mathsf wx,\mathsf wy)-\mathsf l_x(\mathsf ny)-\mathsf l_y(\mathsf nx).
\label{eq:pb-defect-signed-budget}
\end{equation}
Admit exactly those \(\mathsf n\) for which the positive and negative indices \((m_+,m_-)\) of
\(\mathsf G_{\mathsf n}\) obey \(m_+\leq n_+\) and \(m_-\leq n_-\). With
\(\mathcal K_{\mathsf n}=\operatorname{rad}\mathsf G_{\mathsf n}\), choose a complement
\(\mathcal V=\mathcal V_{\mathsf n}\oplus\mathcal K_{\mathsf n}\), an isometric embedding
\(\mathsf I_{\mathsf n}:(\mathcal V_{\mathsf n},\mathsf G_{\mathsf n}|_{\mathcal V_{\mathsf n}})\to
(\mathcal O_1,\Delta_1)\), and a linear map
\(\mathsf Z:\mathcal K_{\mathsf n}\to\mathsf I_{\mathsf n}(\mathcal V_{\mathsf n})^{\perp_{\Delta_1}}\)
with totally isotropic image, \(\Delta_1(\mathsf Zx,\mathsf Zy)=0\) for all \(x,y\). Define
\begin{equation}
\mathsf k(v+k)=\mathsf I_{\mathsf n}v+\mathsf Zk,\qquad
\mathsf J\eta=\mathsf k-\mathsf w+\mathsf n
\quad(v\in\mathcal V_{\mathsf n},\ k\in\mathcal K_{\mathsf n}).
\label{eq:pb-defect-signed-parameterization}
\end{equation}
These are all solutions of \eqref{eq:pb-intrinsic-cocycle-problem}, with arbitrary lifts through
\(\mathsf J\) and ambiguity \(\operatorname{Hom}(\mathcal V,\ker\mathsf J)\); zero spaces are allowed, and
\(\mathsf Z\) may vanish, with rank at most \(\min(n_+-m_+,n_--m_-)\). No Markov hypothesis is used.
The restricted radical cannot be replaced by the ambient one: for
\(\Delta=\begin{psmallmatrix}0&1\\1&0\end{psmallmatrix}\), \(\mathcal O=\operatorname{span}(e_1)\), and
\(\mathsf p(X)=Xe_2\), one has \(\mathcal O_0=\mathcal O\), and \(\mathsf n(X)=\mathrm cXe_1\) gives
\(\mathsf G_{\mathsf n}(X,X)=-2\mathrm cX^2\), so the cocycle forces \(\mathrm c=0\), while for
\(\mathsf p=0\) every \(\mathsf n\) is allowed. \status{ESTABLISHED}

\propositionheading{Smooth solutions on regular strata}{prop:pb-defect-smooth-strata}
Let the preceding data be smooth finite-rank bundles and bundle maps on a paracompact region of the
base. For positivity assume constant ranks of \(\mathsf J\) and \(\mathsf H\) with \(\mathsf H\succeq0\);
for the positive-semidefinite cocycle, constant ranks of \(\Delta\), \(\mathsf J_\Delta\), and
\(\mathsf w_\Delta\). Then the smooth solutions are exactly the smooth contractions or isometric bundle
embeddings of
\cref{eq:pb-defect-contraction-parameterization,eq:pb-defect-isometry-parameterization}, with smooth
lifts and arbitrary smooth kernel-valued maps. For the signed case assume constant ranks of
\(\mathsf J\) and \(\Delta|_{\mathcal O}\), choose a smooth complement \(\mathcal O_1\), and restrict to
smooth \(\mathsf n\) for which \(\mathsf G_{\mathsf n}\) has constant rank and signature; the smooth
solutions are exactly the smooth \(\mathsf I_{\mathsf n},\mathsf Z\) and lifts of
\eqref{eq:pb-defect-signed-parameterization}. Inertia bounds give local smooth \(\mathsf I_{\mathsf n}\)
but not a global isometric embedding into the given \(\mathcal O_1\) bundle, and agreement of the
resulting \(\mathsf k\) on overlaps is part of the parameterization. The classified sets are intrinsic:
passive frame changes carry \(\mathsf D,\Psi_*,\mathsf J,\mathsf p\), their kernels and images, and the
Fisher or signed forms by their tensor laws, and preserve contraction, isometry, inertia, and total
isotropy; auxiliary metrics and splittings change only the parameters of one solution set, no positive
invariant norm on a noncompact adjoint representation is assumed, and nothing extends automatically
across a rank change. \status{ESTABLISHED}

\theoremheading{Exact realization by a target connection}{thm:pb-defect-connection-descent}
Fix the source connection \(\omega\), the principal scale map \(\mathcal P\), and a target reference
connection \(\bar\omega_0\), and let \(\eta_0\) be the descended defect
\(\mathcal P^*\bar\omega_0-\mathrm d\kappa\circ\omega\). An adjoint-valued one-form \(\eta\) is realized
by a target principal connection if and only if
\begin{equation}
\eta=\eta_0+f^*\bar\eta,\qquad
\bar\eta\in\Omega^1(\bar{\mathcal C};\operatorname{Ad}\bar{\mathscr P}).
\label{eq:pb-defect-affine-realization}
\end{equation}
For a surjective smooth submersion \(f\), this holds exactly when \(\delta\eta=\eta-\eta_0\)
annihilates \(\ker Tf\) and satisfies
\begin{equation}
\delta\eta_c(X)=\delta\eta_{c'}(X')
\quad\text{whenever}\quad
f(c)=f(c'),\quad T_cfX=T_{c'}fX',
\label{eq:pb-defect-all-fiber-descent}
\end{equation}
both values lying in the same target adjoint fiber and all components of each fiber included; the
descended \(\bar\eta\) is unique and smooth. The principal-realizable solutions are exactly the
algebraic parameterizations, with all lift freedom retained, intersected with
\eqref{eq:pb-defect-affine-realization}. When \(f\) is a diffeomorphism every smooth adjoint lift is
realizable; for general \(f\) the failure of one lift does not exclude another corrected by a
kernel-valued one-form. No flatness or curvature equation is imposed on the target connection.
\status{ESTABLISHED}

The proofs of the statements of this section and four witnesses for their boundaries, including an
action-rank obstruction to smooth lifts and signed isotropic solutions outside the Markov domain, are in
\Cref{app:base-proofs} (\Cref{prop:pb-defect-witnesses}).
